\documentclass[10pt,letterpaper]{article}

\usepackage[top=1in,bottom=1in,left=1in,right=1in]{geometry}
\usepackage{times}
\usepackage[T1]{fontenc}
\usepackage[utf8]{inputenc}
\usepackage{microtype}
\usepackage{parskip}
\usepackage{titling}
\usepackage{authblk}
\usepackage{hyperref}
\hypersetup{colorlinks=false,hidelinks}

\setlength{\parindent}{0pt}
\setlength{\parskip}{6pt}

% Title formatting to match NECEC template
\pretitle{\begin{center}\fontsize{14}{16}\selectfont\bfseries}
\posttitle{\par\end{center}\vspace{2pt}}
\preauthor{\begin{center}\fontsize{10}{12}\selectfont}
\postauthor{\end{center}}
\predate{}
\postdate{}
\date{}

\title{MamNet-6: Multi-View Fusion, CrossView Transformer, and
Uncertainty-Aware Triage for Breast Cancer Detection on VinDr-Mammo}

\author[1]{Haoming Luo}
\author[1]{Matthew Hamilton}
\author[1]{Oscar Meruvia-Pastor}
\author[1]{Edward Kendall}
\affil[1]{\small Department of Computer Science, Memorial University,
St.\ John's, NL, Canada\\
\texttt{hluo@mun.ca}}

\renewcommand\Affilfont{\small}

\begin{document}

\maketitle
\thispagestyle{empty}
\vspace{-8pt}

\noindent\textbf{Keywords:} breast cancer detection, mammography, multi-view
learning, uncertainty quantification, triage

\vspace{4pt}
\hrule
\vspace{8pt}

\noindent\textbf{Extended Abstract---}Breast cancer is the most commonly
diagnosed cancer worldwide. Population screening with full-field digital
mammography (FFDM) reduces mortality but places a heavy radiologist workload
and substantial inter-reader variability of 10--30 percentage points (Lehman
et al., 2017). Meta-analyses report pooled AUC of 0.87 for deep learning AI
versus 0.81 for radiologists on digital mammography (Dembrower et al., 2020).
Three gaps motivate this work: single-view dominance in published VinDr-Mammo
results, underexplored spatial cross-view attention, and the absence of
uncertainty-aware triage simulation on this dataset.

We developed MamNet-6, a six-stage pipeline evaluated on VinDr-Mammo (Nguyen
et al., 2023), a 20,000-image Vietnamese FFDM benchmark. Each stage extends
the previous: Stage~1 establishes CNN baselines across five backbones;
Stage~2 introduces BreastPairNet, a shared-backbone multi-view network fusing
CC and MLO views across twelve backbone$\times$fusion$\times$resolution
variants; Stage~3 adds a CrossView Transformer with spatial cross-attention
between view feature maps before global average pooling; Stage~3b tests
whether coarse bounding-box annotations help or hurt performance; Stage~4
applies Monte Carlo (MC) Dropout ($N$=20 stochastic passes) for predictive
uncertainty quantification; and Stage~5 assembles a 16-model ensemble with
test-time augmentation (TTA$\times$5). Stage~6 simulates a Verboom-style
hybrid triage policy (Verboom et al., 2021), routing high-confidence AI reads
autonomously and referring uncertain cases to a radiologist. External
zero-shot validation is performed on the RSNA 2022 dataset.

DenseNet-121 achieved AUROC\,=\,0.756 at Stage~1. BreastPairNet improved to
AUROC\,=\,0.826 [95\,\% CI 0.776--0.871] (+9.1\,\%), and the CrossView
Transformer reached AUROC\,=\,0.832 [0.783--0.875] (+10.0\,\%).
Adjacent-stage differences lie within overlapping confidence intervals; the
cumulative Stage~1$\to$5 gain is the most statistically robust comparison.
Bounding-box annotations consistently degraded AUROC by 0.12--0.20,
indicating shortcut learning from coarse localisation labels. The 16-model
ensemble achieved AUROC\,=\,0.846 (+11.8\,\% vs.\ Stage~1) with sensitivity
at 90\,\% specificity (S@90\,\%spec)\,=\,0.657 and positive predictive value
(PPV)\,=\,0.184 at the triage threshold. Hybrid triage at 50\,\% AI reads
preserved 87.9\,\% sensitivity under oracle conditions. Zero-shot transfer to
RSNA 2022 yielded AUROC\,=\,0.601, highlighting a large domain gap between
Vietnamese and North American screening populations.

MamNet-6 contributes a systematic ablation across six design
decisions---backbone, multi-view fusion, spatial cross-attention, uncertainty
quantification, ensemble construction, and clinical triage simulation---on a
single large FFDM benchmark, delivering a cumulative 11.8\,\% AUROC gain over
single-view CNN. Two key negative findings emerge: bounding-box shortcut
learning degrades performance, and domain generalisation to unseen screening
populations remains an open challenge.

\vspace{6pt}
\noindent\textbf{References}

\vspace{2pt}
\noindent Dembrower, K.\ et al.\ (2020). Effect of AI-supported screening on
interval cancers in mammography. \textit{Lancet Oncology}.

\noindent Lehman, C.\ et al.\ (2017). National performance benchmarks for
modern screening digital mammography. \textit{Radiology}.

\noindent Nguyen, H.\ et al.\ (2023). VinDr-Mammo: A large-scale benchmark
dataset for computer-aided detection and diagnosis in full-field digital
mammography. \textit{Scientific Data}.

\noindent Verboom, P.\ et al.\ (2021). Standalone AI for triage in breast
cancer screening. \textit{Radiology}.

\end{document}
